\clearpage\thispagestyle{empty}
\pdfbookmark[0]{Quellen zu Komet (18. Februar 1819)}{cometfeb1819-sources}
\begin{center}\textbf{Quellen zu Komet (18. Februar 1819)}\end{center}
\noindent Diese KI-Übertragung ins heutige Deutsch folgt \textit{Comet} in den \href{https://gdz.sub.uni-goettingen.de/id/PPN236007467}{\textit{Werken}, Band~VI} (Göttingen, 1874), S.~417--419. Der Bericht erschien zuerst unter der Überschrift Göttingen in den \textit{Göttingischen gelehrten Anzeigen}, Stück~28, 18. Februar 1819, S.~\href{https://books.google.com/books?id=CAQEAAAAMAAJ&pg=PA273}{273}--278. Für S.~273--\href{https://books.google.com/books?id=CAQEAAAAMAAJ&pg=PA277}{277} wurde das von Google digitalisierte Exemplar der University of Michigan herangezogen. Der Bericht spricht von Gauß in der dritten Person.
\par\smallskip\noindent In Nicolais vorläufigen parabolischen Elementen gibt der Erstabdruck auf S.~\href{https://books.google.com/books?id=CAQEAAAAMAAJ&pg=PA274}{274} als Logarithmus des kleinsten Abstands 9,52933 an, die Werkausgabe dagegen 0.52933. Auch das Datum 1818 November und die Formulierung „Länge des aufsteigenden Kometen“ beim Vergleich mit dem Kometen von 1805 stehen im Erstabdruck auf S.~\href{https://books.google.com/books?id=CAQEAAAAMAAJ&pg=PA276}{276}. Die Übertragung folgt der Werkausgabe und übernimmt auch die Angabe „etwa 40“ ohne Einheit beim Vergleich der Rektaszensionen. Die historische Fußnote zur wahrscheinlichen Seeberger Zeit steht auch im Erstabdruck auf S.~\href{https://books.google.com/books?id=CAQEAAAAMAAJ&pg=PA275}{275}. KI-Übertragung und KI-Satz.
\clearpage
